\begin{figure}[t]
\centering
\begin{tikzpicture}[>=stealth, node distance=2.5cm, every node/.style={font=\small}]
    % Nodes
    \node[draw, minimum size=1cm, shape=rectangle] (S1) {$s_1$};
    \node[draw, minimum size=1cm, shape=circle, right=of S1] (S2) {$\textcolor{myred}{s_2}$};
    \node[draw, minimum size=1cm, shape=rectangle, right=of S2] (S3) {$\textcolor{myred}{s_3}$};
    \node[draw, minimum size=1cm, shape=circle, right=of S3] (S4) {$\textcolor{myblue}{s_4}$};

    % Self-loops
    \draw[->] (S1) edge[loop above]
        node[above, pos=0.5] {$(s_1, a)$} 
        coordinate[pos=0.5] (S1loop) (S1);
    \filldraw[black] (S1loop) circle (2pt);

    \draw[->] (S3) edge[loop above]
        node[above, pos=0.5] {$\textcolor{myred}{(s_3, a)}$} 
        coordinate[pos=0.5] (S3loop) (S3);
    \filldraw[black] (S3loop) circle (2pt);

    \draw[->] (S4) edge[loop above]
        node[above, pos=0.5] {$\textcolor{myblue}{(s_4, a)}$} 
        coordinate[pos=0.5] (S4loop) (S4);
    \filldraw[black] (S4loop) circle (2pt);

    % Transitions and splits
    \draw[-] (S1) -- node[below] {} ($(S1)!0.5!(S2)$);
    \filldraw[black] ($(S1)!0.5!(S2)$) circle (2pt) 
        node[label={90:$\textcolor{myred}{(s_1, b)}$}] {}; % Labeled split point
    \draw[->] ($(S1)!0.5!(S2)$) -- (S2);
    \draw[->] ($(S1)!0.5!(S2)$) .. controls +(0,-1) and +(0,-1) .. node[below] {} (S1);

    \draw[-] (S2) -- node[below] {} ($(S2)!0.5!(S3)$);
    \filldraw[black] ($(S2)!0.5!(S3)$) circle (2pt) 
        node[label={90:$\textcolor{myred}{(s_2, b)}$}] {};
    \draw[->] ($(S2)!0.5!(S3)$) -- (S3);
    \draw[->] ($(S2)!0.5!(S3)$) .. controls +(0,-2) and +(0,-2) .. node[below] {} (S1);

    \draw[-] (S3) -- node[below] {} ($(S3)!0.5!(S4)$);
    \filldraw[black] ($(S3)!0.5!(S4)$) circle (2pt) 
        node[label={90:$\textcolor{myblue}{\textcolor{myblue}{(s_3, b)}}$}] {};
    \draw[->] ($(S3)!0.5!(S4)$) -- (S4);
    \draw[->] ($(S3)!0.5!(S4)$) .. controls +(0,-3) and +(0,-3) .. node[below] {} (S1);

    % No transition from S4 to S1
\end{tikzpicture}

\caption{An example of a RGN that serves to illustrate the calculation of the weak attractor of the target set $\set{s_1}$ for the square player. In each iteration $C_k$ includes more states, until only $s_1$ and $(s_1, a)$ remain in the weak attractor.}
\label{fig:example_game_02_watr}
\end{figure}